\section{Exclusion of Suzuki and Ree normal cores}
\label{sec:twisted-lie-cores}

This section removes the three Suzuki--Ree series, together with
the Tits group.  The argument allows all normalizing outer
automorphisms and tensor permutations.  In particular, it does
not require a character table for the missing outer extension
of ${}^2F_4(8)$.

\subsection*{D-BI-22a. Complete small-field inputs}

We use $q$ for the actual finite-field size.  Put
\[
 U=2^2.\operatorname{Sz}(8),\qquad A=U.3.
\]
The center of $U$ has order four and is not cyclic.  A nonlinear
irreducible representation of $U$ need not be faithful; retaining
all such representations is essential for covering every possible
quasisimple central quotient.

\begin{lemma}[D-BI-22a: exact small-field bounds]
\label{D-BI-22a}
Every nonscalar element in every nonlinear irreducible representation
of $U$, and in every irreducible representation of $A$ whose
restriction to $U$ is nonlinear and irreducible, satisfies
\[
                         \nu(g)\ge\frac37a\ge\frac14a,
                                                               \tag{BI.85}
\]
where $a$ is the representation degree.  The smallest such degree
is fourteen.

For every nontrivial irreducible character $\chi$ of the three
groups in the following table, and every nonscalar class,
\[
                  |\chi(g)|\le c\chi(1),                  \tag{BI.86}
\]
with the indicated constants and minimum degrees:
\[
\begin{array}{c|c|r}
S&c&\min\chi(1)\\\hline
{}^2G_2(27)&26/703&703\\
{}^2F_4(2)'&3/13&26\\
{}^2F_4(8)&449/32319&64638
\end{array}                                                   \tag{BI.87}
\]
All three constants are at most $1/4$.
\end{lemma}

\begin{proof}
The exact input consists of the complete CTblLib tables named
\texttt{2\^{}2.Sz(8)}, \texttt{2\^{}2.Sz(8).3},
\texttt{R(27)}, \texttt{2F4(2)'}, and \texttt{2F4(8)}.
Their numbers of classes are respectively $35,25,35,22,113$.
The first four have ATLAS source strings.  The last is the
complete table supplied by Gunter Malle, whose provenance is
explicitly recorded in CTblLib, rather than relabelled as an
ATLAS table.  The complete Ree character table in CHEVIE is
also explicitly documented in \cite[\S10.2, numbered p.~131,
before Proposition~10.4]{BI:HHM}; Malle's published unipotent
character computation is \cite{BI:MalleRee}.
The group identifications and character tables are these
published mathematical inputs, not consequences of numerical
orthogonality tests.

The executed program \texttt{twisted\_lie/export\_and\_scan.g}
retains every character, class size, order, prime power map,
source string and the actual $U\to A$ fusion.  It checks full
row orthogonality.  The extension rows are selected by exact
norm-one restriction, giving $34,12,34,21,112$ selected rows,
respectively.  The computation examines all $213$ selected rows
and $15\,471$ nonscalar entries.

For the first two tables it computes all eigenvalue multiplicities
by the exact Fourier formula (BI.81), checks their nonnegative
integrality and degree sums, and finds minimum ratio $3/7$.
For the remaining tables write each retained cyclotomic value as
\[
                 \chi(g)=\sum_j b_j\zeta_m^j,
                 \qquad B(\chi(g))=\sum_j|b_j|.
\]
The coefficients $b_j$ are integers and all basis choices are
retained in the input.  The triangle inequality gives
$|\chi(g)|\le B(\chi(g))$.  Exact rational maximization of
$B(\chi(g))/\chi(1)$ over all nonscalar entries gives precisely
the three constants in (BI.87).  Thus no floating comparison of
cyclotomic values or unproved cancellation is used.

The independent program \texttt{check\_twisted\_lie.py} selects
the rows again using rational normalized-trace restriction norms,
reconstructs all $1330$ Suzuki spectra by Ramanujan Fourier sums,
and recomputes all $14\,141$ remaining coefficient bounds.
It verifies the complete five-table inventory, all minimum degrees,
and the fusion's orders, identity kernel, normal-class unions and
prime power-map compatibility.  Both executions pass.
The complete input is $161615$ bytes with SHA-256
\texttt{10dd43bccf28471f0c8725086666385aba1dc9970fe9d53bd1478e9ad6a1637b}.
These finite exact checks prove the lemma from the stated inputs.
\end{proof}

\subsection*{D-BI-22b. Every single-factor normalizer}

\begin{lemma}[D-BI-22b: universal normalizer bounds]
\label{D-BI-22b}
Let $N$ be a finite quasisimple group acting faithfully and
irreducibly in degree $a$, whose simple quotient is a Suzuki
or Ree group or the Tits group.  Every nonscalar element of
its linear normalizer satisfies
\[
                         \nu(g)\ge\frac{23}{128}a.        \tag{BI.88}
\]
All such factor degrees are at least fourteen.  More specifically:
for $\operatorname{Sz}(8)$ the quarter bound holds; for the
three groups (BI.87) the bound is $3a/16$; and for actual field
size $q\ge32$ the minimum degree is at least $124$.
\end{lemma}

\begin{proof}
The only exceptional Schur multiplier in these three series is
that of $\operatorname{Sz}(8)$; see
\cite[Chapter~5, Table~1, numbered p.~35]{BI:HHM}.
The generic centers are trivial, and the Tits multiplier is trivial.
The small-group universal-cover and outer metadata are also
retained from the official CTblLib extension metadata, interpreted
as in its \S3.7-3 \cite{SPOR:CTblLib}.

For $\operatorname{Sz}(8)$ the universal cover is $U$ above and
the full outer group is cyclic of order three.  Lift the given
irreducible representation to $U$.  An inner normalizer is scalar
times a represented element of $U$, and is covered by (BI.85).
For an outer normalizer, choose the matching generating coset of
the actual extension $A$.  After an inner adjustment its action
on $U$ is the unique lift of the corresponding simple-group
automorphism.  If $t\in A$ represents that coset and $t^3=u_0$,
Schur's lemma gives $g^3=c\rho(u_0)$.  Rescale $g$ by a cube
root of $c^{-1}$.  The conjugation and power relations then give
an irreducible representation of the fixed group $A$ extending
$\rho$, exactly as in D-BI-21c.  Its restriction belongs to the
tested rows.  This proves the quarter bound without assuming a
faithful irreducible representation of the noncyclic center of $U$,
or assuming that every outer automorphism preserves every central
kernel.

For the three groups (BI.87), D-BI-18's complex commutator
argument turns $|\chi(x)|/a\le1/4$ for nonscalar core elements
into
\[
                              \nu(g)\ge3a/16
\]
for every nonscalar normalizer, including every outer coset.
This avoids any need for an outer-extension table for
${}^2F_4(8)$.  The proof of D-BI-18 also handles nonidentity
scalar commutators by their full actual moving rank.

The remaining actual field sizes are at least $32$.
The published ordinary-character bound
\cite[Theorem~4.1(2), p.~141]{BI:DHP} gives
\[
 c(q)=\max\{1/(\sqrt q-1),9/q\}\le9/32.
\]
Indeed $q\ge32$ implies $\sqrt q>5$, so the first term is
less than $1/4<9/32$, and the second is at most $9/32$.
All simple quotients here are nonalternating, as required by
that statement.  The same complex commutator argument gives
$(1-c(q))a/4\ge23a/128$.

For completeness the first-degree formulas, converted from
the square-root field parameter in
\cite[\S\S5.1,5.2,5.5, numbered pp.~9,14--15]{BI:Luebeck},
are
\[
\begin{array}{c|c}
{}^2B_2(q)&(q-1)\sqrt{q/2}\\
{}^2G_2(q)&q^2-q+1\\
{}^2F_4(q)&(q-1)\sqrt{q/2}(q+1)^2(q^2-q+1).
\end{array}                                                   \tag{BI.89}
\]
For fields in this large range their first admissible values are
$q=32,243,32$, respectively, and the corresponding minima are
$124,58807,134090748$.  Each expression is increasing thereafter.
There are no exceptional covers in this range.  At the remaining
small simple groups the minima are those in D-BI-22a, so all
factor degrees are at least fourteen.

Finally $a/4$ and $3a/16$ are both at least $23a/128$.
This proves the uniform assertion and all refinements.
\end{proof}

\subsection*{D-BI-22. The whole Suzuki--Ree source gate}

\begin{theorem}[D-BI-22: Suzuki--Ree core exclusion]
\label{D-BI-22}
Let an inducing block of a finite-cover Weil sixfold source have
finite image $J\subset\operatorname{GL}(V)$.  Suppose $J$ has
an irreducibly acting normal central product of quasisimple
groups whose simple quotients all belong to
\[
 {}^2B_2(2^{2m+1}),\quad{}^2G_2(3^{2m+1}),\quad
 {}^2F_4(2^{2m+1})\quad(m\ge1),\qquad{}^2F_4(2)'.
\]
Then this source has period rank at most eight.

The statement allows mixed factor types, all compatible outer
automorphisms, arbitrary scalar phases, and tensor permutations.
It requires neither primitive or tensor-indecomposable $J$,
essential self-duality, equality of inducing and original fields,
nor a nonsplit-polarization assumption.
\end{theorem}

\begin{proof}
Write the normal central product as $N=N_1\cdots N_t$.
Its irreducible representation has the tensor form
$V=V_1\otimes\cdots\otimes V_t$, with faithful irreducible
factor representations and degrees $a_i\ge14$.
Put $e=\prod_i a_i$.
The factors $N_i$ are intrinsic: each is the derived subgroup
of the inverse image of its simple direct factor in $N/Z(N)$.
Hence $J$ permutes them.  Their full matrix algebras
$\operatorname{End}(V_i)\otimes1$, supplied by irreducibility,
are permuted as well.  Thus a normalizing operator is a tensor
product of factor intertwiners followed by a permutation of
isomorphic factor spaces.

For a factor-preserving nonscalar operator at least one factor
intertwiner is nonscalar.  Fixing the eigenvalues in all other
factors shows that its largest eigenspace fraction cannot exceed
that in this factor.  D-BI-22b therefore gives
$\nu(g)\ge23e/128$.

If an operator permutes factors nontrivially, choose one cycle
of length $l\ge2$ on factors of dimension $a$.
After tensor changes of basis the action is a cyclic word rotation
with a diagonal twist at one place.  Such a diagonalization is
possible because the operator has finite order up to scalar.
On each orbit of basis words the weighted cyclic action has at
most one vector for each fixed eigenvalue.  There are $a$ constant
words; every other orbit has size at least two.  Its largest
eigenspace therefore has dimension at most
\[
                  a+(a^l-a)/2\le3a^l/4.
\]
Tensoring with all other cycles preserves this fraction, giving
$\nu(g)\ge e/4\ge23e/128$.  This is the same elementary
weighted-word argument as D-SPOR-03, now with no restriction on
the component types.  Genuine scalar inertia moves all $e$
dimensions.

Assume for contradiction that the actual period rank is nine.
The proved intermediate-cover factorization of D-BI-01 applies
to this inducing block, with arbitrary inducing field.  If its
intermediate base has genus $h$ and $R$ genuine branch points,
the multiplicity-six Euler identity is
\[
                    D=6+2e(1-h),\qquad R\le D.
\]
At $h=1$ the actual variation factors through a finite Hurwitz
cover of dimension at most six.  At $h\ge2$ the displayed $D$
is negative since $e\ge14$.  Hence $h=0$.  The same actual
period factorization then forces $R\ge12$, while $D=2e+6$.
Applying the uniform normalizer bound to every genuine inertia
gives the necessary integer and continuous inequalities
\[
  12\left\lceil\frac{23e}{128}\right\rceil\le2e+6,
                 \qquad e\le\frac{192}{5}.                \tag{BI.90}
\]
This uses an upper bound for the actual period map, not an
inference of dominance from parameter counting.

If $t\ge2$, then $e\ge14^2=196$, contrary to (BI.90).
For a single factor of field size at least $32$, one has
$e\ge124$, again impossible.  For a single
$\operatorname{Sz}(8)$ factor the stronger quarter bound gives
$3e\le2e+6$, so $e\le6<14$.
For each remaining single small factor, the stronger $3/16$
bound gives $9e/4\le2e+6$, so $e\le24$, whereas its degree
is at least $26$ by (BI.87).  Every case is impossible.
The target period domain has dimension nine, so its rank is at
most eight, as asserted.
\end{proof}

\paragraph{Boundary of the result.}
The nonsimple group ${}^2G_2(3)$ is not added to the displayed
simple-family list.  Its derived simple group is
$\operatorname{PSL}_2(8)$ and belongs to the separately treated
$A_1$ case.  No claim about the other classical or exceptional
Lie families is inferred from the Suzuki--Ree calculation.
